% ===== PRÉAMBULE CANONIQUE — à coller VERBATIM en tête de CHAQUE .tex =====
\documentclass[11pt,a4paper]{article}
\usepackage[utf8]{inputenc}\usepackage[T1]{fontenc}\usepackage{lmodern}
\usepackage[french]{babel}
\usepackage{amsmath,amssymb,mathtools}\usepackage{icomma}
\usepackage[version=4]{mhchem}          % \ce{...} : équations & formules
\usepackage{siunitx}\sisetup{locale=FR} % \qty{}{}, \si{}, \ang{}
\usepackage{chemfig}                    % \chemfig{...} : structures développées
\usepackage{xcolor}\usepackage{colortbl}
\usepackage{tikz}\usepackage{pgfplots}\pgfplotsset{compat=1.18}
\usepackage[most]{tcolorbox}\usepackage{enumitem}\usepackage{array,booktabs}
\usepackage[a4paper,margin=2.2cm]{geometry}
\usetikzlibrary{arrows.meta,calc,angles,quotes,positioning,patterns,backgrounds,shapes.geometric}
\definecolor{primary}{RGB}{222,49,99}\definecolor{primarydark}{RGB}{150,22,55}
\definecolor{defcol}{RGB}{0,114,178}\definecolor{methcol}{RGB}{52,92,148}
\definecolor{excol}{RGB}{0,135,90}\definecolor{attncol}{RGB}{200,40,55}
\tcbset{boxbase/.style={enhanced,breakable,boxrule=0.9pt,arc=2.5pt,
  left=3mm,right=3mm,top=2.3mm,bottom=2.3mm,fonttitle=\bfseries,coltitle=white}}
% --- boîtes TUTORIEL (noms EXACTS, ne pas renommer) ---
\newtcolorbox{cle}[1][]{boxbase,colback=defcol!6,colframe=defcol,
  title={\textsc{À retenir}\ifx\relax#1\relax\relax\else\textmd{ : \itshape #1}\fi}}
\newtcolorbox{methode}[1][]{boxbase,colback=methcol!7,colframe=methcol,sharp corners,
  title={\textsc{Méthode}\ifx\relax#1\relax\relax\else\textmd{ --- \itshape #1}\fi}}
\newtcolorbox{exemplebox}[1][]{boxbase,colback=excol!5,colframe=excol,
  title={\textsc{Exemple}\ifx\relax#1\relax\relax\else\textmd{ : \itshape #1}\fi}}
\newtcolorbox{piege}[1][]{boxbase,colback=attncol!6,colframe=attncol,
  title={\textsc{Piège}\ifx\relax#1\relax\relax\else\textmd{ : \itshape #1}\fi}}
\newtcolorbox{remarque}[1][]{boxbase,colback=black!4,colframe=black!45,
  title={\textsc{Remarque}\ifx\relax#1\relax\relax\else\textmd{ : \itshape #1}\fi}}
% --- boîtes EXERCICE / CORRIGÉ (noms EXACTS, auto-comptées) ---
\newtcolorbox[auto counter]{exo}[1][]{boxbase,colback=primary!4,colframe=primary,
  title={\textsc{Exercice}~\thetcbcounter\ifx\relax#1\relax\relax\else\textmd{ --- \itshape #1}\fi}}
\newtcolorbox[auto counter]{corrige}[1][]{boxbase,colback=excol!4,colframe=excol,
  title={\textsc{Corrigé}~\thetcbcounter\ifx\relax#1\relax\relax\else\textmd{ --- \itshape #1}\fi}}
\newcommand{\R}{\mathbb{R}}\newcommand{\N}{\mathbb{N}}\newcommand{\Z}{\mathbb{Z}}\newcommand{\ds}{\displaystyle}
% =========================================================================
\begin{document}
\sisetup{per-mode=symbol}
\begin{corrige}[précipitation lors d'un mélange]
\begin{enumerate}[label=\alph*)]
  \item Le volume total est \qty{100}{\milli\litre} : chaque concentration est divisée par 2.
    \[ [\ce{Pb^{2+}}] = 0{,}020\times\dfrac{50}{100} = \qty{0.010}{\mole\per\litre}, \qquad
       [\ce{I-}] = 0{,}020\times\dfrac{50}{100} = \qty{0.010}{\mole\per\litre}. \]
  \item $Q = [\ce{Pb^{2+}}][\ce{I-}]^{2} = (0{,}010)(0{,}010)^{2} = \num{1.0e-6}$. Comme
        $Q=\num{1.0e-6} > K_{\mathrm{ps}}=\num{7.1e-9}$, la solution est sursaturée : \textbf{\ce{PbI2}
        précipite} (précipité jaune).
\end{enumerate}
\end{corrige}

\begin{corrige}[effet d'ion commun sur \ce{AgCl}]
\begin{enumerate}[label=\alph*)]
  \item $s_0 = \sqrt{K_{\mathrm{ps}}} = \sqrt{\num{1.8e-10}} = \qty{1.34e-5}{\mole\per\litre}$.
  \item $s' = [\ce{Ag+}] = \dfrac{K_{\mathrm{ps}}}{[\ce{Cl-}]} = \dfrac{\num{1.8e-10}}{0{,}10}
           = \qty{1.8e-9}{\mole\per\litre}$.
  \item Facteur $= \dfrac{s_0}{s'} = \dfrac{\num{1.34e-5}}{\num{1.8e-9}} \approx 7400$ : la solubilité
        est divisée par environ \num{7000}.
\end{enumerate}
\end{corrige}

\begin{corrige}[ion commun avec un anion doublement présent]
\begin{enumerate}[label=\alph*)]
  \item $[\ce{Pb^{2+}}]=s'$ et $[\ce{Br-}]$ est imposée, donc
        $K_{\mathrm{ps}}=[\ce{Pb^{2+}}][\ce{Br-}]^{2}=s'\,[\ce{Br-}]^{2}$, d'où
        $s'=\dfrac{K_{\mathrm{ps}}}{[\ce{Br-}]^{2}}$.
  \item $s' = \dfrac{\num{9.2e-6}}{(0{,}50)^{2}} = \dfrac{\num{9.2e-6}}{0{,}25}
           = \qty{3.7e-5}{\mole\per\litre}$ (contre $s_0=\qty{1.3e-2}{\mole\per\litre}$ dans l'eau pure :
        chute d'un facteur $\sim 360$).
\end{enumerate}
\end{corrige}

\begin{corrige}[concentration seuil de précipitation]
\begin{enumerate}[label=\alph*)]
  \item La précipitation débute quand $Q=K_{\mathrm{ps}}$, soit $[\ce{Ag+}][\ce{Cl-}]=K_{\mathrm{ps}}$ :
    \[ [\ce{Cl-}]_{\min} = \dfrac{K_{\mathrm{ps}}}{[\ce{Ag+}]} = \dfrac{\num{1.8e-10}}{\num{1.0e-3}}
       = \qty{1.8e-7}{\mole\per\litre}. \]
  \item En dessous de $\qty{1.8e-7}{\mole\per\litre}$, $Q<K_{\mathrm{ps}}$ : la solution est non
        saturée, \textbf{aucun} précipité ne se forme (tout reste dissous).
\end{enumerate}
\end{corrige}

\begin{corrige}[y a-t-il précipitation de \ce{CaCO3} ?]
On compare $Q=[\ce{Ca^{2+}}][\ce{CO3^{2-}}]$ à $K_{\mathrm{ps}}=\num{3.3e-9}$.
\begin{enumerate}[label=\alph*)]
  \item $Q = (\num{1.0e-3})(\num{1.0e-5}) = \num{1.0e-8} > K_{\mathrm{ps}}$ : \textbf{précipitation}.
  \item $Q = (\num{1.0e-4})(\num{1.0e-5}) = \num{1.0e-9} < K_{\mathrm{ps}}$ : \textbf{pas} de précipité.
\end{enumerate}
\end{corrige}
\end{document}
